\section{The secondary scale and positive quadratic drift}
\label{sec:quadratic-regimes}

When the normalized ratios of \(q_0\) and \(q_1\) to the scale set by
\(q_2\) are small, we use the secondary scales \(s,r\) of
\cref{sec:standard-tools}.
For \(q_2>0\), we compare the projected boundary flows near their first
roots and continue the half-orbits to the fixed output windows.
We use curvature away from the symmetric limit and the exterior
contraction in root-distance coordinates near it.
The case \(q_2<0\) is treated in
\cref{sec:direct-negative-annulus-buffer,sec:direct-negative-root-regimes}:
we first estimate the annular passage and the negative branch, then
treat the positive branch in each coefficient region.
For fixed parameters, the early, annular or root, and large-exit ranges
form overlapping subintervals of \(\mathcal H_p\). Their estimates give
crossing directions or critical-point curvature for \(\mathcal L\),
which are combined by \cref{cor:ordered-cover-zero-count}.
Thus the coefficient region selects a passage argument, while the
ordered height cover combines the estimates within that argument.
\Cref{fig:secondary-cover} shows how these regions cover the normalized
coefficients of both signs of \(q_2\).

\begin{figure}[H]
\centering
\begin{tikzpicture}[x=4.6cm,y=3.2cm,>=Stealth,font=\small,
    bar/.style={line width=2.6pt,gray!55},
    lab/.style={font=\footnotesize,inner sep=1.5pt,text=black}]
  % (a) negative quadratic coefficient
  \begin{scope}
    \draw[->] (0,0)--(1.2,0) node[right] {$a$};
    \draw[->] (0,0)--(0,1.62) node[above] {$b$};
    \draw[very thick] (0,1)--(1,1)--(1,0);
    \draw (0,1)--(-.03,1) node[left] {$1$};
    \draw (1,0)--(1,-.03) node[below] {$1$};
    \draw (.25,1)--(.25,.96) node[below,lab] {$\tfrac14$};
    \draw[bar] (.01,1.09)--(.14,1.09) node[lab,right] {weak root};
    \draw[bar] (.07,1.2)--(.22,1.2) node[lab,right] {separated roots};
    \draw[bar] (.19,1.31)--(.31,1.31) node[lab,right] {fold};
    \draw[bar] (.28,1.42)--(1,1.42);
    \node[lab,above] at (.64,1.44) {buffered};
    \draw[bar] (1.08,0)--(1.08,1);
    \node[lab,right] at (1.1,.5) {buffered};
    \node at (.5,-.25) {(a) \(q_2<0\)};
  \end{scope}
  % (b) positive quadratic coefficient
  \begin{scope}[xshift=7.4cm]
    \draw[->] (0,0)--(1.2,0) node[right] {$a$};
    \draw[->] (0,0)--(0,1.62) node[above] {$b$};
    \draw[very thick] (0,1)--(1,1)--(1,0);
    \draw (0,1)--(-.03,1) node[left] {$1$};
    \draw (1,0)--(1,-.03) node[below] {$1$};
    \draw[bar] (.01,1.09)--(.16,1.09) node[lab,right] {weak root};
    \draw[bar] (.08,1.2)--(1,1.2);
    \node[lab,above] at (.6,1.22) {curvature};
    \draw[bar] (1.08,.2)--(1.08,1);
    \draw[bar] (1.16,0)--(1.16,.4);
    \node[lab,right] at (1.1,.72) {curvature};
    \node[lab,right] at (1.18,.2) {small \(b\)};
    \node at (.5,-.25) {(b) \(q_2>0\)};
  \end{scope}
\end{tikzpicture}
\caption{Schematic of the compact coefficient covers on the secondary
charts \eqref{eq:direct-negative-secondary-scales} and
\eqref{eq:direct-positive-secondary-chart}; thresholds are not to scale.
The normalized coefficients lie on the curve \(\max\{\sqrt a,b\}=1\), and
the bars show the overlapping regions. (a) For \(q_2<0\): the buffered
region \(b^2-4a\le-\nu\) (\cref{prop:direct-buffered-negative}), the fold
\(|b^2-4a|\le2\nu\) (\cref{prop:direct-negative-fold}), separated roots
\(b^2-4a\ge\nu\), \(a\ge a_0^-\) (\cref{prop:direct-negative-separated}), and
the weak root \(a\le a_-\)
(\cref{prop:direct-negative-weak-root}). (b) For \(q_2>0\): curvature for
\(a\ge a_0^+\), \(b\ge b_{*,+}/2\), root-distance coordinates for small \(b\), and
the weak root \(a\le a_+\) (\cref{prop:direct-positive-quadratic}).
Near \(a=0\), the cover combines the weak-root regions, which require
\(s/a\le\kappa_\pm\), with the contact strip of
\cref{prop:direct-bounded-contact}.}
\label{fig:secondary-cover}
\end{figure}

